\subsection{TrajGaze Encoder}
\label{sec:encoder}

The TrajGaze encoder maps a past trajectory sequence
$(G_{1:T_p}, H^L_{1:T_p}, H^R_{1:T_p})$ and an optional text query $q$
to a per-patch importance score $\mathbf{s} \in [0,1]^{P}$
over the $P{=}196$ spatial positions of the visual backbone.
The encoder is composed of four sequential stages: per-modality tokenization,
intra-frame cross-stream attention, inter-frame temporal reasoning, and
visual--trajectory fusion.

\paragraph{Temporal Input Construction.}
Given a raw clip of arbitrary length, we uniformly subsample $T{=}32$ frames.
For each frame we record the gaze position $g_t \in \mathbb{R}^2$,
gaze speed $\dot{g}_t \in \mathbb{R}$ (frame-to-frame displacement),
and left/right hand positions $h_t^{L,R} \in \mathbb{R}^2$ together with
their frame-to-frame velocity $\dot{h}_t^{L,R} \in \mathbb{R}^2$ and a binary
visibility flag $m_t^{L,R} \in \{0,1\}$.
The 32 frames are divided into a \emph{past} prefix of length $T_p$ and a
\emph{future} suffix of length $T_f = T - T_p$.
During Stage~1 pretraining $T_p / T$ is drawn uniformly from $[0.4, 0.6]$ each
iteration to prevent the encoder from overfitting to a fixed context length.
Only the past prefix is given to the encoder; the future is used solely for supervision.

\paragraph{Per-Modality Tokenization.}
\label{sec:tokenizer}
For each past frame $t \in \{1, \dots, T_p\}$ we construct four tokens
via independent learned linear projections:
\begin{align}
    z_t^{g}   &= \mathrm{LN}\!\left(W_g\,[g_t;\dot{g}_t] + b_g\right), \\
    z_t^{L}   &= \mathrm{LN}\!\left(W_L\,[h_t^L;\dot{h}_t^L;m_t^L] + b_L\right), \\
    z_t^{R}   &= \mathrm{LN}\!\left(W_R\,[h_t^R;\dot{h}_t^R;m_t^R] + b_R\right), \\
    z_t^{\phi} &= \mathrm{LN}\!\left(W_\phi\,\phi(g_t, h_t^L, h_t^R) + b_\phi\right),
\end{align}
where $[\cdot;\cdot]$ denotes vector concatenation, $\mathrm{LN}$ is layer normalization,
and all tokens are projected to dimension $d{=}128$.
The interaction feature $\phi \in \mathbb{R}^{12}$ captures pairwise geometric and
kinematic relations: Euclidean distances from gaze to each hand
$(d^L, d^R) \in \mathbb{R}^{2 \times 2}$, relative velocities
$(v^L_{\mathrm{rel}}, v^R_{\mathrm{rel}}) \in \mathbb{R}^{2 \times 2}$,
a scalar gaze--hand convergence score, and a scalar temporal lead-lag indicator.
When a modality is undetected (visibility flag off, or gaze tracking failure),
the corresponding linear projection is bypassed and replaced by a dedicated
learned \emph{missing} embedding $e_{\mathrm{miss}} \in \mathbb{R}^d$,
allowing the model to distinguish absent from zero-valued readings.

\paragraph{Intra-Frame Attention.}
The four tokens of each frame are passed jointly through a 2-layer, 4-head
transformer encoder~\cite{vaswani2017attention} to model cross-stream
interactions within the same timestep:
\begin{equation}
    \bigl[z_t^g,\, z_t^L,\, z_t^R,\, z_t^\phi\bigr]
    \;\leftarrow\;
    \mathrm{TF}_{L_1}\!\left(\bigl[z_t^g,\, z_t^L,\, z_t^R,\, z_t^\phi\bigr]\right),
    \quad t = 1,\dots,T_p.
\end{equation}
All $T_p$ frames are processed in parallel ($B T_p$ independent sequences of length 4).
This stage allows the gaze token to be informed by the current hand configuration
before any temporal reasoning occurs.

\paragraph{Inter-Frame Temporal Transformer.}
The $T_p \times 4$ intra-frame tokens are projected to $D{=}256$ dimensions,
flattened into a single sequence of length $T_p \cdot 4$, and augmented with
sinusoidal positional encodings~\cite{vaswani2017attention}:
\begin{equation}
    \mathbf{x} = \mathrm{PE}\!\left(
        W_{\mathrm{proj}}\,\mathbf{z}_{1:T_p}
    \right) \in \mathbb{R}^{(T_p \cdot 4) \times D}.
\end{equation}
Because each frame contributes exactly four consecutive positions, the sinusoidal
encoding implicitly encodes both the temporal index and the within-frame token type.
A 6-layer, 8-head transformer encoder then applies full self-attention across the
entire sequence:
\begin{equation}
    \mathbf{x} \leftarrow \mathrm{TF}_{L_2}(\mathbf{x})
    \in \mathbb{R}^{(T_p \cdot 4) \times D},
\end{equation}
enabling each token at any timestep and modality to attend to all others.
This stage captures long-range temporal patterns such as gaze anticipation,
hand motion build-up, and recurrent interaction signatures.

\paragraph{Query-Conditioned Feature Modulation.}
A FiLM~\cite{perez2018film} layer conditions the trajectory context on the
text query embedding $e_q \in \mathbb{R}^{d_q}$
(produced by the frozen CLIP~\cite{radford2021clip} text encoder):
\begin{equation}
    \tilde{\mathbf{x}} = \mathbf{x} \odot \bigl(1 + \gamma(e_q)\bigr) + \beta(e_q),
\end{equation}
where $\gamma, \beta : \mathbb{R}^{d_q} \to \mathbb{R}^D$ are learned linear projections.
During Stage~1 pretraining the query is absent and $e_q{=}\mathbf{0}$ (FiLM becomes
an identity); at inference the query shifts which visual patches the trajectory
context attends to, enabling task-specific patch selection.

\paragraph{Visual--Trajectory Fusion.}
\label{sec:vtfusion}
A frozen DINOv2-ViT/S backbone~\cite{oquab2023dinov2} extracts spatial patch
features $\mathbf{V} \in \mathbb{R}^{P \times D_v}$ from a single representative frame.
We perform multi-head cross-attention with the trajectory context as queries and
visual patches as keys and values:
\begin{equation}
    A^h = \mathrm{softmax}\!\left(
        \frac{(\tilde{\mathbf{x}} W_Q^h)\,(\mathbf{V} W_K^h)^\top}{\sqrt{d_h}}
    \right) \in \mathbb{R}^{(T_p \cdot 4) \times P}, \quad h=1,\dots,H,
\end{equation}
with $H{=}8$ heads and head dimension $d_h = D/H$.
Intuitively, each trajectory token asks: ``which visual patch explains my current
state?''  The enriched trajectory context
$\tilde{\mathbf{x}} + W_O\!\left(\bigoplus_h A^h \mathbf{V} W_V^h\right)$
is passed to the prediction decoders.

The per-patch importance score is derived from the aggregate attention:
\begin{equation}
    \bar{A}[p] = \frac{1}{H} \sum_{h=1}^{H}
                  \frac{1}{T_p \cdot 4} \sum_{i=1}^{T_p \cdot 4} A^h[i, p],
\end{equation}
which measures how consistently each patch was queried across all trajectory
tokens and all heads.  A lightweight score head
$f_s : \mathbb{R}^D \to [0,1]$ then refines this into the final importance score:
\begin{equation}
    \mathbf{s}[p] = \sigma\!\left(
        W_s \,\mathrm{LN}\!\left(\bar{A}[p] \cdot v_p\right) + b_s
    \right) \in [0,1],
\end{equation}
where $v_p$ denotes the visual feature of patch $p$ and $\sigma$ is the sigmoid function.

\paragraph{Stage~1 Pretraining.}
\label{sec:stage1}
Before joint fine-tuning with the VLM, the TrajGaze encoder is pretrained on a
trajectory-prediction task to acquire a meaningful notion of gaze-relevant patches
without requiring question--answer supervision.
Given the past trajectory prefix $(1{:}T_p)$, the encoder is trained to predict:
(i) future gaze positions $\hat{g}_{T_p+1:T}$ via a 3-layer cross-attention
trajectory decoder, and
(ii) future per-patch importance scores $\hat{\mathbf{s}}_{T_p+1:T} \in [0,1]^{P}$
via a score decoder.
The ground-truth importance score at frame $t$ is computed as a temporally
smoothed Gaussian centered at the observed gaze position, multiplied by a
hand-proximity weight:
\begin{equation}
    I(p,\,t) =
        \underbrace{G(p,\,t)}_{\text{gaze term}}
        \cdot
        \underbrace{H(p,\,t)}_{\text{hand term}}
        \cdot \Phi_t \cdot \Psi_t,
\end{equation}
where $G(p,t) = \exp(-\|c_p - g_t\|^2 / 2\sigma_g^2)$,
$H(p,t)$ is the analogous hand-proximity Gaussian, and
$\Phi_t, \Psi_t$ are the gaze--hand convergence and lead-lag scalars.
The Stage~1 training loss is:
\begin{equation}
    \mathcal{L}_{\mathrm{S1}} =
        \underbrace{\mathcal{L}_{\mathrm{traj}}}_{\text{MSE on future gaze}}
        + \underbrace{\mathcal{L}_{\mathrm{score}}}_{\text{BCE on future patch scores}}
        + \underbrace{\mathcal{L}_{\mathrm{attn}}}_{\text{MSE: patch\_scores vs.\ mean}(I_{\mathrm{future}})}.
\end{equation}
The third term directly supervises the encoder's attention weights
($\mathbf{s}$ from~\S\ref{sec:vtfusion}) to align with the temporal average of
future ground-truth scores, ensuring that what the encoder \emph{attends to}
while processing the past matches what will actually be important in the future.
Pretraining uses clips from the EgoExo4D-learn and HoloAssist domains (246 unique clips,
same domain as Stage~3 fine-tuning) to minimize the distribution gap.
